% SPDX-License-Identifier: CC-BY-4.0
% Copyright (C) 2026 Fabian Januszewski, Abhijit (the Langlandscape team)
\section{A sign in the literature: \texorpdfstring{$\JL(U_q)$}{JL(U\_q)} for \texorpdfstring{$q\mid N^-$}{q | N-}}\label{sec:hsieh}
Our work relies throughout on Hsieh's paper \cite{Hs}, whose results on Gorenstein Hecke algebras and congruence ideals of
definite quaternionic Hida families enter stages~04 and~06. While fixing normalisations we found one sentence where, as
we read it, a sign should be corrected. We record it here briefly; the details are in the technical note \cite{LLhsieh}
(\texttt{mathematics/notes/hsieh-sign.pdf}, 4 pp.). Fabian intends to contact the author after the Mathathon.

\paragraph{The statement} \stcited. In \cite[\S4.5, p.~34]{Hs} (arXiv:1705.02717v5; page numbers of the published version
not checked), the Jacquet--Langlands homomorphism $\JL:\T(N,\mathrm{I})\to\T^D(N,\mathrm{I})$ is stated to satisfy $\JL(U_q)=(-1)U_q$ for
$q\mid N^-$.

\paragraph{The correction} (\stpmc{} for trivial central character at $q$, i.e.\ \stproved{} modulo \stcited{} inputs; for an
unramified non-trivial central character at $q$ the same argument is a \stsketch; \cite[Prop.~2.1]{LLhsieh}). With the
normalisations of \cite{Hs} the relation is $\JL(U_q)=+U_q$; equivalently, the $\varpi^D_q$-eigenvalue of the primitive
Jacquet--Langlands lift $f^D$ is $\lambda^D_f(U_q)=a(q,f)$. In our conventions: $U_q\mapsto+\Pi_q$. All other values of $\JL$
in \cite[p.~34]{Hs} are unchanged. The argument is local: $\pi_q=(\mu\circ\det)\otimes\St$ corresponds to $\pi'_q=\mu\circ\Nrd$,
and both eigenvalues equal $\mu(q)$ times the same positive factor. Read literally, the printed sign would make the module
$e\mathbf{S}^D(N,\mathrm{I})[\lambda^D_f]$ of \cite[Thm~4.5]{Hs} zero \cite[Lemma~2.3]{LLhsieh}; with the corrected sign it is the intended line.

\paragraph{Evidence} \stcomp. T05-1 (\res{s05/T05-1.json}): on all $144$ eigenlines of $S_k(\Gamma_0(q))^{\rm new}$,
$q\in\{11,23,37,43\}$, $k\in\{2,4,6,8\}$, the Brandt eigenvalue of $\Pi_q$ equals $+a_q$; both signs of the eigenvalue occur
(70 and 74 lines), so the test separates the two relations. Independent checks: the weight-2 geometric description in
\cite{Ri90}, the level-11 case ($a_{11}(11a)=+1$, $\Pi_{11}$ the identity on the class set), and the convention of \cite{CH15}
(as cited in the note).

\paragraph{Consequences.} As far as we can see the correction concerns one normalisation sentence; we found no effect on the
theorems of \cite{Hs} (Theorem~4.5 and its proof, the theta elements, the congruence ideals and the main theorems use the
local representations directly or are intrinsic to $\T^D$) \cite[\S4]{LLhsieh}; this is our reading, not a complete audit of
\cite{Hs}. In this project we use $U_q\mapsto+\Pi_q$ (decision FD2, Fabian, 15 November 2026), and no statement of \cite{Hs}
involving $U_q$ for $q\mid N^-$ is used with the printed sign. A plausible origin is the Atkin--Lehner relation $w_q=-U_q$ on
$q$-new forms of weight~2 or the sign in the local character relation (\stconj).
